% Problem record op_25342dbb8d64e728. This TeX file is derived from
% database/problems_json/op_25342dbb8d64e728.json by scripts/sync-tex.mjs; edit the JSON record
% and rerun the script rather than editing this file.
\section{High-dimensional crossover of optimized port-based teleportation}
\label{sec:op_25342dbb8d64e728}

\paragraph{Problem.}
What is the high-dimensional crossover fidelity of optimized deterministic port-based teleportation?

Let $F_d^*(N)$ be the maximum entanglement fidelity of teleporting an unknown $d$-dimensional state with $N$ ports, optimizing both the resource state and Alice's measurement while Bob only selects a port. Use

\begin{equation}
F_e(\mathcal T)=\operatorname{Tr}\!\left[\Phi_d(\operatorname{id}\otimes\mathcal T)(\Phi_d)\right],
\qquad
\Phi_d=\frac1d\sum_{a,b=1}^d|aa\rangle\langle bb|.
\label{eq:2534-1}
\end{equation}

For each fixed $c>0$, define

\begin{equation}
N_d=\max\{1,\lfloor cd^2\rfloor\},
\qquad
\varphi(c)=\lim_{d\to\infty}F_d^*(N_d).
\label{eq:2534-2}
\end{equation}

Does the limit in Eq.~\eqref{eq:2534-2} exist for every $c>0$, and if so, what is $\varphi(c)$? Equation~\eqref{eq:2534-1} fixes the fidelity convention.

\subsection*{Status}

Unsolved

\subsection*{Source}

The question is explicitly posed or retained as open in the cited primary literature \sourcecite{ref:2534-christandl21}{Christandl21}\sourcecite{ref:2534-yoshida26}{Yoshida26}.  The statement is rewritten here to make its hypotheses and success criterion self-contained.

\subsection*{Progress}

\begin{itemize}
  \item Finite-resource optimization is already characterized by a teleportation-matrix eigenvalue. In particular,
  \begin{equation}
F_d^*(N)=\frac{N}{d^2}\qquad(1\leq N\leq d).
\label{eq:2534-4}
\end{equation}
  This exact regime does not reach $N\sim c d^2$ for fixed positive $c$. \sourcecite{ref:2534-mozrzymas18}{Mozrzymas18}

The displayed definitions, constraints, and target bounds are recorded in Eqs.~\eqref{eq:2534-4}.

  \item The optimized exact probabilistic protocol has success probability
  \begin{equation}
p_d^*(N)=\frac{N}{N+d^2-1}.
\label{eq:2534-5}
\end{equation}
  Converting failure into an arbitrary port output gives a deterministic protocol with fidelity at least $p_d^*(N)$. Consequently,
  \begin{equation}
\liminf_{d\to\infty}F_d^*(N_d)\geq\frac{c}{1+c}.
\label{eq:2534-6}
\end{equation}
  Section 8 explicitly identifies the fixed-ratio limit as an open direction. \sourcecite{ref:2534-christandl21}{Christandl21}

The displayed definitions, constraints, and target bounds are recorded in Eqs.~\eqref{eq:2534-5}, \eqref{eq:2534-6}.

  \item An elementary no-signalling consequence of the protocol model gives
  \begin{equation}
F_d^*(N)\leq\min\left\{1,\frac{N}{d^2}\right\}.
\label{eq:2534-7}
\end{equation}
  Indeed, for a fixed port $B_i$, its subnormalized reference–output state in branch $i$ is bounded above by the unconditional state $\mathbb I_R/d\otimes\rho_{B_i}$. Its overlap with $\Phi_d$ is therefore at most $1/d^2$; summing over $N$ selected-port branches proves the bound. Combining this derivation with the preceding lower bound yields the useful, not claimed sharp, bracket
  \begin{equation}
\frac{c}{1+c}
  \leq\liminf_{d\to\infty}F_d^*(N_d)
  \leq\limsup_{d\to\infty}F_d^*(N_d)
  \leq\min\{1,c\}.
\label{eq:2534-8}
\end{equation}
  This paragraph supplies a direct derivation rather than attributing the bracket to a new theorem of the cited paper. \sourcecite{ref:2534-christandl21}{Christandl21}

The displayed definitions, constraints, and target bounds are recorded in Eqs.~\eqref{eq:2534-7}, \eqref{eq:2534-8}.

  \item The 2026 correspondence relates the problem to unitary estimation with $N_d-1$ queries, but its order estimate $1-F_d^*(N)=\Theta(d^4N^{-2})$ does not determine $\varphi(c)$. In particular, a fixed-$d$ expansion cannot be substituted into this joint limit without uniform error control. \sourcecite{ref:2534-yoshida26}{Yoshida26}
\end{itemize}

\subsection*{References}

\begin{enumerate}[label={},leftmargin=4.8em,labelsep=0.5em]
  \footnotesize
  \item[\textup{[Mozrzymas18]}]\label{ref:2534-mozrzymas18}
  M. Mozrzymas, M. Studziński, S. Strelchuk, and M. Horodecki, "Optimal Port-based Teleportation," \emph{New Journal of Physics} 20, 053006 (2018). \href{https://doi.org/10.1088/1367-2630/aab8e7}{doi:10.1088/1367-2630/aab8e7}; \href{https://arxiv.org/abs/1707.08456}{arXiv:1707.08456}.

  \item[\textup{[Christandl21]}]\label{ref:2534-christandl21}
  M. Christandl, F. Leditzky, C. Majenz, G. Smith, F. Speelman, and M. Walter, "Asymptotic performance of port-based teleportation," \emph{Communications in Mathematical Physics} 381, 379–451 (2021). \href{https://doi.org/10.1007/s00220-020-03884-0}{doi:10.1007/s00220-020-03884-0}; \href{https://arxiv.org/abs/1809.10751}{arXiv:1809.10751}.

  \item[\textup{[Yoshida26]}]\label{ref:2534-yoshida26}
  S. Yoshida, Y. Koizumi, M. Studziński, M. T. Quintino, and M. Murao, "One-to-One Correspondence between Deterministic Port-Based Teleportation and Unitary Estimation," \emph{IEEE Transactions on Information Theory} 72, 2358–2377 (2026). \href{https://doi.org/10.1109/TIT.2026.3658543}{doi:10.1109/TIT.2026.3658543}; \href{https://arxiv.org/abs/2408.11902}{arXiv:2408.11902}.
\end{enumerate}

\subsection*{Comment}

A literature check through 15 September 2026 located no determination of this crossover function or general proof of the displayed limit. This is distinct from the fixed-dimensional asymptotic-coefficient problem: the input dimension grows together with the resource, so dimension-dependent remainders matter. The result would quantify high-dimensional port requirements rather than extrapolating a fixed-dimensional approximation.

\subsection*{Field}

Quantum Communication

\subsection*{Topic}

Channel simulation; One-shot and finite-blocklength bounds

\subsection*{ID}

\texttt{op\_25342dbb8d64e728}
